\section{Aplicaciones de Meta-RL}

\subsection{Manipulación robótica}

Meta-RL tiene en la robótica un campo de aplicación natural: un robot necesita adaptarse con rapidez a objetos distintos, a objetivos de agarre distintos, a parámetros dinámicos distintos, etcétera.

\subsection{In-Context Learning en los modelos de lenguaje de gran tamaño}

\begin{knowledgebox}{In-Context Learning como Meta-Learning}
El in-context learning de los modelos de lenguaje de gran tamaño (resolver una tarea nueva a partir de los ejemplos incluidos en el prompt) puede verse como una forma de context-based meta-learning. Durante el preentrenamiento, el LLM ha visto una gran cantidad de «tareas» y ha aprendido a inferir la tarea a partir del contexto y a generar la salida correspondiente. Esto sigue directamente la idea de RL$^2$.
\end{knowledgebox}

\subsection{Resumen de la sección}

Las ideas de Meta-RL ya se han extendido a varios campos, como la robótica y el procesamiento de lenguaje natural; el in-context learning puede verse como el surgimiento natural del meta-learning en el preentrenamiento a gran escala.
